\documentclass[12pt,a4paper]{article}

% ── Kódovanie a jazyk ─────────────────────────────────────────────────────────
\usepackage[utf8]{inputenc}
\usepackage[T1]{fontenc}
\usepackage[slovak]{babel}

% ── Písmo ─────────────────────────────────────────────────────────────────────
\usepackage{lmodern}
\usepackage{microtype}

% ── Matematika ────────────────────────────────────────────────────────────────
\usepackage{amsmath, amssymb, amsthm}

% ── Grafika a tabuľky ─────────────────────────────────────────────────────────
\usepackage{graphicx}
\usepackage{booktabs}
\usepackage{array}
\usepackage{float}

% ── Formátovanie ──────────────────────────────────────────────────────────────
\usepackage{geometry}
\geometry{top=2.5cm, bottom=2.5cm, left=3cm, right=2.5cm}
\usepackage{setspace}
\onehalfspacing
\usepackage{parskip}
\setlength{\parindent}{0pt}

% ── Zdrojový kód ──────────────────────────────────────────────────────────────
\usepackage{listings}
\usepackage{xcolor}

\lstdefinelanguage{JavaScript}{
  keywords={function,var,let,const,return,if,else,for,while,new,this,class,true,false,null},
  keywordstyle=\bfseries,
  comment=[l]{//},
  morecomment=[s]{/*}{*/},
  string=[b]",
  morestring=[b]',
  sensitive=true
}

\lstset{
  language=JavaScript,
  basicstyle=\ttfamily\small,
  frame=single,
  framerule=0.4pt,
  rulecolor=\color{black!30},
  backgroundcolor=\color{black!4},
  breaklines=true,
  breakatwhitespace=true,
  numbers=left,
  numberstyle=\tiny\color{black!40},
  numbersep=6pt,
  xleftmargin=1.5em,
  captionpos=b,
  showstringspaces=false
}

% ── Obsah, hypertextové prepojenia ────────────────────────────────────────────
\usepackage{hyperref}
\hypersetup{colorlinks=true,linkcolor=black,urlcolor=black,citecolor=black}

% ── Pomocné definície ─────────────────────────────────────────────────────────
\newcommand{\norm}[1]{\left\lVert#1\right\rVert}
\newcommand{\T}{\mathsf{T}}
\newcommand{\gitrepo}{https://github.com/tvoje-meno/svm-projekt}

% =============================================================================
\begin{document}

% ── Titulná strana ────────────────────────────────────────────────────────────
\begin{titlepage}
  \centering
  \vspace*{2cm}
  {\LARGE\bfseries Dokumentácia projektu\par}
  \vspace{0.6cm}
  {\Huge\bfseries Metóda podporných vektorov\\[0.3em]
   \large (Support Vector Machine -- SVM)\par}
  \vspace{1.2cm}
  {\large Predmet: Strojové učenie\par}
  \vspace{0.4cm}
  {\large TUKE FEI - ÚUI - 2025/2026\par}
  \vspace{0.8cm}
  {\large Autor: Samuel Štefán\par}
  \vspace{0.3cm}
  {\large Stránka: \url{https://samuelstefan01.github.io/ml-svm-website}\par}
  \vspace{0.3cm}
  {\small Repozitár: \url{https://github.com/SamuelStefan01/ml-svm-website}\par}
  \vfill
  \rule{\linewidth}{0.5pt}\\[0.4em]
  {\small Interaktívna webová aplikácia implementovaná v~HTML/CSS/JavaScript}
\end{titlepage}

% ── Obsah ─────────────────────────────────────────────────────────────────────
\tableofcontents
\newpage

% =============================================================================
\section{Úvod a ciele projektu}
% =============================================================================

Cieľom projektu je interaktívna webová aplikácia, ktorá komplexne prezentuje
algoritmus \textbf{Support Vector Machine (SVM)} -- metódu podporných vektorov.
Aplikácia pokrýva teoretický základ, matematickú formuláciu, vizualizáciu
rozhodovacej hranice, vyhodnotenie metrík efektívnosti a porovnanie s~inými
klasifikačnými algoritmami.

Projekt je implementovaný ako~jediný HTML súbor (\texttt{index.html})
bez externých závislostí okrem knižnice Chart.js (verzia 4.4.1, načítaná z~CDN).
Algoritmus SMO (Sequential Minimal Optimization) je napísaný od základov
v~JavaScripte bez použitia externých knižníc strojového učenia.

\subsection*{Štruktúra aplikácie}

Aplikácia je rozdelená do piatich záložiek:
\begin{enumerate}
  \item \textbf{Teória SVM} -- matematická formulácia, geometria marginu,
        stratégie viactriednej klasifikácie (OvO/OvR)
  \item \textbf{Interaktívna demo} -- klikateľné plátno, výber jadra
        a~parametrov, vizualizácia podporných vektorov a~váhového vektora
  \item \textbf{Jadrové funkcie} -- porovnanie lineárneho, RBF
        a~polynomiálneho jadra na rovnakom datasete
  \item \textbf{Datasety a metriky} -- $k$-násobná krížová validácia,
        ROC krivka, mriežkové vyhľadávanie, krivka učenia
  \item \textbf{Porovnanie modelov} -- SVM vs.\ $k$-NN vs.\ logistická
        regresia vs.\ rozhodovací strom
\end{enumerate}

% =============================================================================
\section{Teoretický základ}
% =============================================================================

\subsection{Definícia SVM}

SVM je algoritmus učenia s~učiteľom (\textit{supervised learning}), ktorý hľadá
\textbf{optimálnu nadrovinu} (hyperplane) oddeľujúcu dáta do dvoch tried
so~\textbf{maximálnym marginom}. Margin je vzdialenosť medzi rozhodovacou
hranicou a~najbližšími dátovými bodmi z~každej triedy. Tieto kľúčové body
sa nazývajú \textbf{podporné vektory} (support vectors) a~sú to jediné
trénovacie vzorky, ktoré model potrebuje na predikciu.

\subsection{Geometria marginu}

\begin{figure}[H]
  \centering
  \includegraphics[width=0.82\textwidth]{geometria_marginu.png}
  \caption{Geometria SVM: rozhodovacia hranica, marginový pás a~podporné
           vektory. Šírka marginu je $2/\norm{\mathbf{w}}$.}
  \label{fig:geometria}
\end{figure}

Rozhodovacia hranica je definovaná rovnicou:
\begin{equation}
  \mathbf{w} \cdot \mathbf{x} + b = 0,
\end{equation}
pričom body spĺňajúce $\mathbf{w} \cdot \mathbf{x} + b = +1$ tvoria kladnú
marginovú nadrovinu a~body spĺňajúce $\mathbf{w} \cdot \mathbf{x} + b = -1$
tvoria zápornú marginovú nadrovinu. Šírka marginu je teda $2/\norm{\mathbf{w}}$
(viď obrázok \ref{fig:geometria}).

\subsection{Optimalizačný problém -- primárna forma}

Pre \textbf{lineárne separovateľné dáta} (hard-margin SVM):
\begin{equation}
  \min_{\mathbf{w},\,b} \;\frac{1}{2}\norm{\mathbf{w}}^2
  \quad \text{s.t.} \quad
  y_i(\mathbf{w} \cdot \mathbf{x}_i + b) \geq 1, \quad \forall\, i.
\end{equation}

Pre \textbf{lineárne neseparovateľné dáta} (soft-margin SVM) sa zavádza
premenná ohybnosti (slack variable) $\xi_i \geq 0$ a~regularizačný
parameter $C > 0$:
\begin{equation}
  \min_{\mathbf{w},\,b,\,\boldsymbol{\xi}} \;\frac{1}{2}\norm{\mathbf{w}}^2
  + C\sum_{i=1}^{n} \xi_i
  \quad \text{s.t.} \quad
  y_i(\mathbf{w} \cdot \mathbf{x}_i + b) \geq 1 - \xi_i,\quad \xi_i \geq 0.
\end{equation}

Parameter $C$ riadi kompromis medzi veľkosťou marginu a~počtom chybných
klasifikácií: vysoké $C$ zodpovedá tvrdému marginu (riziko pretrénovania),
nízke $C$ zodpovedá mäkkému marginu (riziko podtrénovania).

\subsection{Duálna forma a trik s jadrom}

Pomocou Lagrangeových multiplikátorov $\alpha_i \geq 0$ sa problém prevádza
do~duálnej formy:
\begin{equation}
  \max_{\boldsymbol{\alpha}} \;\sum_{i=1}^{n}\alpha_i
  - \frac{1}{2}\sum_{i=1}^{n}\sum_{j=1}^{n}
    \alpha_i\alpha_j y_i y_j K(\mathbf{x}_i, \mathbf{x}_j)
  \quad \text{s.t.} \quad
  0 \leq \alpha_i \leq C,\quad \sum_{i=1}^{n}\alpha_i y_i = 0.
\end{equation}

V~duálnej forme sa vstupné vektory vyskytujú iba cez skalárny súčin
$\mathbf{x}_i \cdot \mathbf{x}_j$, ktorý možno nahradiť \textbf{jadrovou
funkciou} $K(\mathbf{x}_i, \mathbf{x}_j)$. Toto je \emph{trik s~jadrom}
(kernel trick): implicitne mapuje dáta do~vyššieho dimenzionálneho priestoru
$\varphi(\mathbf{x})$ bez explicitného výpočtu transformácie, keďže
$K(\mathbf{x}_i, \mathbf{x}_j) = \varphi(\mathbf{x}_i)\cdot\varphi(\mathbf{x}_j)$.

Predikčná funkcia natrénovaného modelu:
\begin{equation}
  f(\mathbf{x}) = \sum_{i \in \mathrm{SV}} \alpha_i y_i
  K(\mathbf{x}_i, \mathbf{x}) + b,
\end{equation}
kde sumácia prebieha iba cez podporné vektory (body s~$\alpha_i > 0$).

\subsection{Jadrové funkcie}

\begin{table}[H]
  \centering
  \caption{Prehľad jadrových funkcií implementovaných v~aplikácii}
  \begin{tabular}{llp{5.8cm}}
    \toprule
    \textbf{Jadro} & \textbf{Vzorec} & \textbf{Typické použitie} \\
    \midrule
    Lineárne
      & $K(\mathbf{x},\mathbf{y}) = \mathbf{x}^{\T}\mathbf{y}$
      & Textová klasifikácia, riedke vysokodimenzionálne dáta \\[4pt]
    RBF (Gaussovo)
      & $K(\mathbf{x},\mathbf{y}) = e^{-\gamma\norm{\mathbf{x}-\mathbf{y}}^2}$
      & Väčšina nelineárnych problémov; predvolené jadro \\[4pt]
    Polynomiálne
      & $K(\mathbf{x},\mathbf{y}) = (\gamma\,\mathbf{x}^{\T}\mathbf{y}+c)^d$
      & Spracovanie obrazu, polynomiálne vzťahy príznakov \\
    \bottomrule
  \end{tabular}
  \label{tab:jadra}
\end{table}

Parameter $\gamma$ v~RBF jadre riadi \emph{dosah vplyvu} každého
trénovacieho bodu: vysoké $\gamma$ zodpovedá úzkemu dosahu (riziko
pretrénovania), nízke $\gamma$ zodpovedá širšiemu dosahu (riziko
podtrénovania).

\subsection{Viactriedna klasifikácia: OvO a OvR}

Základný SVM je binárny klasifikátor. Pre $K > 2$ tried sa používajú dve
stratégie rozkladu:

\begin{description}
  \item[One-vs-Rest (OvR)] Pre každú z~$K$ tried sa natrénuje jeden SVM
    (trieda $k$ oproti všetkým ostatným). Celkovo $K$ modelov. Predikuje
    trieda s~najvyšším skóre $f_k(\mathbf{x})$.
  \item[One-vs-One (OvO)] Pre každý pár tried sa natrénuje jeden SVM.
    Celkovo $K(K-1)/2$ modelov. Predikuje trieda s~najviac hlasmi
    (väčšinové hlasovanie).
\end{description}

Knižnica \texttt{scikit-learn} používa predvolene OvR pre lineárne jadro
a~OvO pre ostatné jadrá. Pri $K = 10$ (napr.\ MNIST) vyžaduje OvO
45~modelov, zatiaľ čo OvR iba 10.

% =============================================================================
\section{Implementácia}
% =============================================================================

\subsection{Algoritmus SMO}

Na~riešenie kvadratickej optimalizačnej úlohy je implementovaný algoritmus
\textbf{Sequential Minimal Optimization (SMO)} podľa Platta (1998).
SMO rozkladá veľký QP problém na sériu minimálnych dvojrozmerných
podproblémov, ktoré majú analytické riešenie.

Kľúčové kroky implementácie:
\begin{enumerate}
  \item Predvýpočet celej matice jadra $K_{ij}$ pre zrýchlenie tréningu
  \item Heuristika výberu druhého pivota: maximalizácia $|E_i - E_j|$
        (Plattova heuristika -- výber najväčšieho kroku)
  \item Analytická aktualizácia $\alpha_i$, $\alpha_j$ s orezaním
        na $[L, H]$
  \item Aktualizácia prahu $b$ podľa podmienok KKT
  \item Konvergencia: zastavenie po niekoľkých priebehoch bez zmeny
\end{enumerate}

\begin{lstlisting}[caption={Jadro SMO aktualizácie -- heuristika a analytická aktualizácia}]
// Heuristika vyberu j: maximalizuj |Ei - Ej|
let best = -1, bestDiff = 0;
for (let k = 0; k < n; k++) {
  if (k === i) continue;
  const diff = Math.abs(Ecache[i] - Ecache[k]);
  if (diff > bestDiff) { bestDiff = diff; best = k; }
}

// Analytická aktualizácia alpha_j
const eta = 2*K[i][j] - K[i][i] - K[j][j];
if (eta >= 0) continue;
this.alpha[j] -= y[j] * (Ecache[i] - Ej) / eta;
this.alpha[j] = Math.max(L, Math.min(H, this.alpha[j]));

// Aktualizácia prahu b
const b1 = this.b - Ei
  - y[i]*(this.alpha[i]-ai0)*K[i][i]
  - y[j]*(this.alpha[j]-aj0)*K[i][j];
this.b = (this.alpha[i] > 0 && this.alpha[i] < C) ? b1 : ...;
\end{lstlisting}

Pre lineárne jadro sa po tréningu explicitne vypočíta váhový vektor:
\begin{equation}
  \mathbf{w} = \sum_{i=1}^{n} \alpha_i y_i \mathbf{x}_i.
\end{equation}
Tento vektor je vizualizovaný v~aplikácii ako šípka v~2D priestore
spolu so skutočnou šírkou marginu $2/\norm{\mathbf{w}}$.

\subsection{Implementované klasifikátory pre porovnanie}

Okrem SVM sú pre účely porovnania implementované od základov:

\begin{itemize}
  \item \textbf{$k$-najbližší susedia ($k$-NN)}: predikcia na základe
    väčšinového hlasovania $k$ najbližších trénovacích bodov
    (euklidovská vzdialenosť, $k = 5$).
  \item \textbf{Logistická regresia}: gradientový zostup minimalizujúci
    krížovú entropiu, 150 epoch, rýchlosť učenia $\eta = 0.1$.
  \item \textbf{Rozhodovací strom}: rekurzívne delenie pomocou Giniho
    nečistoty, maximálna hĺbka $d = 5$.
\end{itemize}

\subsection{Datasety}

\begin{table}[H]
  \centering
  \caption{Datasety použité v~aplikácii}
  \begin{tabular}{lllll}
    \toprule
    \textbf{Dataset} & \textbf{Zdroj} & \textbf{Vzorky}
      & \textbf{Príznaky} & \textbf{Triedy} \\
    \midrule
    Iris (setosa vs.\ versicolor)
      & UCI ML Repository & 75 + 75 & 4 & 2 \\
    Breast Cancer (benígne vs.\ malígne)
      & UCI ML Repository & 36 + 28 & 4 & 2 \\
    Mesiačiky (syntetický)
      & generovaný & variabilný & 2 & 2 \\
    Kruhy (syntetický)
      & generovaný & variabilný & 2 & 2 \\
    XOR (syntetický)
      & generovaný & variabilný & 2 & 2 \\
    \bottomrule
  \end{tabular}
  \label{tab:datasety}
\end{table}

Reálne datasety sú normalizované na interval $[0, 1]$ pomocou min-max
škálovania samostatne pre každý príznak.

% =============================================================================
\section{Metriky efektívnosti}
% =============================================================================

\subsection{Základné metriky klasifikácie}

Z~konfúznej matice (TP, FP, FN, TN) sú vypočítavané nasledujúce metriky:

\begin{align}
  \text{Presnosť (Accuracy)} &= \frac{TP + TN}{TP + FP + FN + TN}, \\[6pt]
  \text{Precíznosť (Precision)} &= \frac{TP}{TP + FP}, \\[6pt]
  \text{Úplnosť (Recall)} &= \frac{TP}{TP + FN}, \\[6pt]
  \text{F1-skóre} &= \frac{2 \cdot \text{Precíznosť} \cdot \text{Úplnosť}}
    {\text{Precíznosť} + \text{Úplnosť}}, \\[6pt]
  \text{Špecificita (Specificity)} &= \frac{TN}{TN + FP}.
\end{align}

Konfúzna matica je v~aplikácii zobrazovaná s~farebnými poliami:
pravé pozitívy (TP) žltou, pravé negatívy (TN) oranžovou,
nepravé pozitívy (FP) červenou a~nepravé negatívy (FN) ružovou.

\subsection{$k$-násobná krížová validácia}

Na~spoľahlivé vyhodnotenie výkonnosti modelu sa používa
\textbf{$k$-násobná krížová validácia} ($k$-fold cross-validation).
Dataset sa rozdelí na~$k$ podmnožín (foldov); model sa $k$-krát
natrénuje na~$k-1$ foldoch a~testuje na~zostávajúcom folde.
Výsledná metrika je priemer a~smerodajná odchýlka:
\begin{equation}
  \bar{A} = \frac{1}{k}\sum_{f=1}^{k} A_f,
  \qquad
  \sigma = \sqrt{\frac{1}{k}\sum_{f=1}^{k}(A_f - \bar{A})^2}.
\end{equation}
Aplikácia podporuje $k \in \{3, 5, 10\}$.

\subsection{ROC krivka a AUC}

ROC krivka (Receiver Operating Characteristic) znázorňuje závislosť
miery pravých pozitív (TPR) od~miery nepravých pozitív (FPR) pri~rôznych
klasifikačných prahoch. Plocha pod~krivkou (AUC) sa odhaduje
lichobežníkovou metódou:
\begin{equation}
  \mathrm{AUC} \approx \sum_{i=1}^{m}
  \bigl(\mathrm{FPR}_i - \mathrm{FPR}_{i-1}\bigr)
  \cdot \frac{\mathrm{TPR}_i + \mathrm{TPR}_{i-1}}{2}.
\end{equation}
Hodnota AUC = 1 zodpovedá dokonalému klasifikátoru, AUC = 0.5
zodpovedá náhodnému uhádnutiu.

\subsection{Mriežkové vyhľadávanie (Grid Search)}

Mriežkové vyhľadávanie systematicky vyskúša všetky kombinácie hodnôt
parametrov $C$ a~$\gamma$ zo~zadanej mriežky:
\begin{equation}
  C \in \{0.1,\; 0.5,\; 1,\; 5,\; 10,\; 50\},
  \qquad
  \gamma \in \{0.05,\; 0.1,\; 0.5,\; 1,\; 5,\; 10\}.
\end{equation}
Pre~každú kombináciu sa vykonáva 3-násobná krížová validácia. Výsledky
sú zobrazené ako~tepelná mapa (heatmap): tmavá farba = nízka presnosť,
svetlá zlatá = vysoká presnosť. Optimálna kombinácia je zvýraznená rámčekom.

\subsection{Krivka učenia (Learning Curve)}

Krivka učenia zobrazuje závislosť trénovacej a~testovacej presnosti
od~veľkosti trénovacej množiny (10\,\% až 90\,\% dostupných dát).
Umožňuje diagnostikovať:
\begin{itemize}
  \item \textbf{Pretrénovanie (overfitting)}: veľká medzera medzi
    trénovacou a~testovacou krivkou (medzera $> 20\,\%$)
  \item \textbf{Podtrénovanie (underfitting)}: obe krivky konvergujú
    k~nízkej hodnote ($< 70\,\%$)
  \item \textbf{Dobrá generalizácia}: krivky sa priblížia k~sebe
    pri~vysokej testovacej presnosti
\end{itemize}
Aplikácia automaticky vyhodnotí medzeru a~zobrazí textové odporúčanie
(napr.\ znížiť $C$ pri pretrénovaní).

% =============================================================================
\section{Popis funkcionality aplikácie}
% =============================================================================

\subsection{Interaktívna vizualizácia}

Záložka \emph{Interaktívna demo} umožňuje:
\begin{itemize}
  \item kliknutím na~plátno (canvas $500 \times 500$ px) pridávať body
    dvoch tried,
  \item výber jadrovej funkcie (lineárne, RBF, polynomiálne),
  \item nastavenie parametrov $C \in [10^{-2},\; 10^{3}]$
    a~$\gamma \in [10^{-2},\; 10^{2}]$ pomocou posuvníkov,
  \item načítanie prednastavených datasetov (lineárny, XOR, kruhy,
    mesiačiky),
  \item vizualizáciu tepelnej mapy predikčného skóre, marginových
    priamok a~zvýraznenie podporných vektorov s~aureolou.
\end{itemize}

Pre~\textbf{lineárne jadro} sa po tréningu zobrazuje sekcia
\emph{Váhový vektor $\mathbf{w}$} obsahujúca:
\begin{itemize}
  \item smerový SVG diagram vektora $\mathbf{w}$ v~2D priestore
    (šípka + prerušovaná kolmica znázorňujúca smer hranice),
  \item relatívne veľkosti zložiek $w_x$, $w_y$ ako~animované
    progress bary,
  \item normu $\norm{\mathbf{w}}$ a~šírku marginu $2/\norm{\mathbf{w}}$.
\end{itemize}

\subsection{Metriky zobrazované po tréningu}

\begin{table}[H]
  \centering
  \caption{Metriky zobrazované po tréningu v~interaktívnej demo}
  \begin{tabular}{ll}
    \toprule
    \textbf{Metrika} & \textbf{Popis} \\
    \midrule
    Accuracy
      & Podiel správne klasifikovaných trénovacích bodov \\
    Support vectors
      & Počet podporných vektorov (body s~$\alpha_i > 0$) \\
    Margin $2/\norm{\mathbf{w}}$
      & Šírka marginu (len pre lineárne jadro) \\
    SMO iterácie
      & Počet iterácií optimalizačného algoritmu \\
    \bottomrule
  \end{tabular}
\end{table}

\subsection{Porovnanie algoritmov}

Záložka \emph{Porovnanie modelov} vyhodnocuje 4 algoritmy pomocou
4-násobnej krížovej validácie na~syntetickom datasete podľa výberu.
Výsledky sú prezentované ako:
\begin{itemize}
  \item vodorovný stĺpcový graf presnosti s~farebne odlíšenými stĺpcami,
  \item vizualizácia rozhodovacích hraníc pre~všetky štyri algoritmy
    vedľa seba (tepelná mapa na plátne $190 \times 190$ px),
  \item súhrnná tabuľka s~presnosťou, smerodajnou odchýlkou, F1-skóre
    a~časom tréningu v~milisekundách.
\end{itemize}

% =============================================================================
\section{Technické riešenie}
% =============================================================================

\subsection{Technológie}

\begin{table}[H]
  \centering
  \caption{Použité technológie}
  \begin{tabular}{lll}
    \toprule
    \textbf{Technológia} & \textbf{Verzia} & \textbf{Účel} \\
    \midrule
    HTML5             & --      & Štruktúra dokumentu, Canvas API \\
    CSS3              & --      & Štýlovanie, animácie, CSS premenné \\
    JavaScript (ES6+) & --      & Logika, SMO, ML algoritmy \\
    Chart.js          & 4.4.1   & Grafy (ROC, krivka učenia, bar chart) \\
    SVG               & --      & Geometrický diagram, váhový vektor \\
    \bottomrule
  \end{tabular}
\end{table}

\subsection{Architektúra kódu}

Celý kód je uložený v jedinom súbore bez externých závislostí
na knižniciach strojového učenia. Hlavné triedy a funkcie:

\begin{description}
  \item[\texttt{class SVM}] Implementácia SVM s~voliteľným jadrom.
    Metódy: \texttt{K()} -- jadrová funkcia, \texttt{train()} -- SMO
    s~heuristickým výberom pivota, \texttt{predict()} -- predikcia.
  \item[\texttt{generateDataset(type, n)}] Generátor syntetických
    datasetov (mesiačiky, kruhy, XOR, lineárny).
  \item[\texttt{kFoldCV(X, y, k, kernel, C, gamma)}] Univerzálna
    $k$-násobná krížová validácia; vracia priemerné Accuracy,
    smerodajnú odchýlku a~priemerné F1-skóre.
  \item[\texttt{computeMetrics(yTrue, yPred)}] Výpočet TP, FP, FN, TN
    a~odvodených metrík vrátane F1-skóre a~špecificity.
  \item[\texttt{runGridSearch()}] Mriežkové vyhľadávanie parametrov
    $C \times \gamma$ s~vykreslením tepelnej mapy.
  \item[\texttt{runLearningCurve()}] Krivka učenia s~automatickou
    interpretáciou výsledku (overfitting / underfitting / OK).
  \item[\texttt{trainSVM()}] Obsluha interaktívneho plátna,
    vizualizácia váhového vektora pomocou SVG.
\end{description}

% =============================================================================
\section{Výsledky a zhrnutie}
% =============================================================================

Aplikácia demonštruje, že SVM je výkonný algoritmus s~jasnou geometrickou
interpretáciou a~flexibilnou jadrovovou funkciou. Kľúčové zistenia
z~experimentov:

\begin{itemize}
  \item \textbf{Lineárne jadro} dosahuje vysokú presnosť na~lineárne
    separovateľných datasetoch (Iris: $\approx 95\,\%$), avšak zlyháva
    na~nelineárnych dátach (kruhy, XOR).
  \item \textbf{RBF jadro} je najuniverzálnejšie -- dosahuje $> 95\,\%$
    presnosť na~mesiačikoch aj kruhoch pri~správnom nastavení $\gamma$.
  \item \textbf{Polynomiálne jadro} ($d = 3$) funguje dobre pre~XOR
    a~iné polynomiálne vzťahy, ale je citlivé na~voľbu parametrov.
  \item \textbf{Mriežkové vyhľadávanie} jednoznačne ukazuje, že voľba
    optimálnych parametrov $C$ a~$\gamma$ má zásadný vplyv na~výkonnosť.
  \item Na~testovaných datasetoch SVM s~RBF jadrom konzistentne
    prekonáva logistickú regresiu na~nelineárnych problémoch a~dosahuje
    porovnateľné výsledky s~$k$-NN pri~výrazne menšom počte uložených
    trénovacích vzoriek.
\end{itemize}

% =============================================================================
\section*{Referencie}
\addcontentsline{toc}{section}{Referencie}
% =============================================================================

\begin{thebibliography}{9}

\bibitem{cortes1995}
  C.~Cortes, V.~Vapnik,
  \textit{Support-Vector Networks}.
  Machine Learning, 20(3):273--297, 1995.

\bibitem{platt1998}
  J.~C.~Platt,
  \textit{Sequential Minimal Optimization: A Fast Algorithm for Training
  Support Vector Machines}.
  Microsoft Research Technical Report MSR-TR-98-14, 1998.

\bibitem{scholkopf2002}
  B.~Schölkopf, A.~J.~Smola,
  \textit{Learning with Kernels: Support Vector Machines, Regularization,
  Optimization, and Beyond}.
  MIT Press, 2002.

\bibitem{uci}
  D.~Dua, C.~Graff,
  \textit{UCI Machine Learning Repository}.
  University of California, Irvine, 2019.
  \url{https://archive.ics.uci.edu}

\end{thebibliography}

\end{document}